\documentclass[../main.tex]{subfiles}
\begin{document}
%==========================================TIÊU ĐỀ TẬP========================================
\begin{table}[h]
    \begin{adjustwidth}{-1cm}{}
        \begin{tabular}{>{\centering\arraybackslash}p{6cm}|>{\raggedright\arraybackslash}p{12.5cm}}
            \multirow{5}{6cm}
            % Cột bên trái: ÉPISODE và số tập
            {\\[-40pt]
                \flaregothic\fontsize{20pt}{20pt}\selectfont \centering
                \color{eptype}PERSONNALISÉ\color{black}\\[12pt]
                \vnmdisp\fontsize{36pt}{36pt}\selectfont
                \color{epnum}NĂM
                \\[18pt]
                % \flaregothic\fontsize{14pt}{14pt}\selectfont
                % \color{epattr}BIENVENUE, \uline{\color{Banpire} Mel} !
                % \flaregothic\fontsize{18pt}{18pt}\selectfont
                % \color{epattr}ÉPISODE PERDU
                \color{black}
            }
             &
            % Dòng thứ nhất: tên tập tiếng Nhật
            {\fotskip\fontsize{18pt}{18pt}\selectfont るしあの60\ruby{秒}{びょ}\ruby{作}{さく}\ruby{戦}{せん}
            }                                                                                   \\[6pt]
            % Dòng thứ hai: tên tập tiếng Anh
             & {\cafeteria\fontsize{18pt}{18pt}\selectfont Baseball Madness}                                                                                                      \\[4pt]
            % Dòng thứ ba: tên tập tiếng Việt
             & {\vnmsans\fontsize{18pt}{18pt}\selectfont Cuộc tập dượt bóng chày}                                                                                                 \\[8pt]
            % Dòng thứ tư: Nhân vật xuất hiện
             & {\sen{\fontsize{14pt}{14pt}\selectfont Track used: Crisscross (\ruby{経}{けい}\ruby{天}{てん}\ruby{緯}{い}\ruby{地}{ち}　\ruby{機}{き}\ruby{略}{りゃく}\ruby{縦}{じゅう}\ruby{横}{おう})}}
            \\[18pt]
            % % Dòng cuối cùng: Thời gian diễn ra của tập (thường sẽ là ngày phát hành)
            % & {\continuum\fontsize{16pt}{16pt}\selectfont Ngày phát hành: 12/09/2021}\\[18pt]
        \end{tabular}
    \end{adjustwidth}
\end{table}
%=========================================TRUYỆN CHÍNH========================================
\color{black}

Giữa tháng Mười, gió thu đã bắt đầu có vị se lạnh, những cơn gió mỏng manh lùa qua bên trong văn phòng \hll. Mặt trời chậm rãi trườn về phía tây, để lại những vệt sáng khẽ xiên tràn lên mặt thảm, vi vu những làn gió trung thu mát lạnh.

Rushia đứng ở giữa phòng, cầm một quả bóng chày cũ kỹ trong tay. Trên người cô là một chiếc áo khoác mỏng màu xám tro, tóc buộc búi cao như mọi ngày, mái tóc xanh lơ lắc lư theo nhịp gió. Cô gái nhỏ trông có vẻ chẳng hợp với môn thể thao này chút nào, ấy vậy mà ánh mắt lại nghiêm túc đến kỳ lạ.

``Bởi vì\dots'' cô nói, giọng nhỏ xíu như đang độc thoại, ``nếu không làm gì đó, mọi người sẽ quên lãng mình mất\dots''

Câu nói ấy không có người chứng kiến. Nhưng chỉ vài phút sau, Pekora xuất hiện từ phía cửa ra vào, tay cầm một cây gậy nhựa màu hồng sáng, vừa đi vừa nhai cà rốt. Bộ đồ thể thao của cô hơi nhàu, rõ ràng là mặc đại vào cho hợp bối cảnh hơn là vì yêu thể thao.

``Bà chắc là muốn chơi trò này thật hả, peko?'' cô hỏi, gậy gác lên vai, mắt nheo nheo nhìn nàng Pháp sư như đnag nghi ngờ ai đó vừa rủ mình chơi đánh cờ dưới mưa.

Rushia gật đầu: ``Ừm. Một trận duy nhất thôi-nanodesu.''

``Không có luật lệ kỳ quặc nào chứ, peko? Không đánh bom, không cho gọi hồn, không cho gọi fan đến cổ vũ, chứ?''

``Ừm,'' nàng Pháp sư gật đầu. ``Chỉ là một trận bóng chày bình thường thôi.''

Ngay lúc đó, từ phía cầu thang bên hông tòa nhà, Korone xuất hiện, khoác ba lô, đầu vẫn còn đeo tai nghe. Cô ngáp dài, vươn vai: ``Nghe nói có vận động nhẹ để giãn gân cốt? Cho chị chơi với!''

``Thế là đủ ba người rồi!'' Rushia nói, có vẻ hơi vui vì kế hoạch nhỏ của mình đang thành hình.

``Khoan đã,'' Pekora ngẩng đầu, ``ai làm trọng tài đây, peko?''

Câu hỏi ấy như mũi tên bay vào cửa kính phòng kỹ thuật, nơi Kuon đang vừa uống cà phê vừa kiểm tra bản dựng trang ứng dụng hôm qua.

Và dĩ nhiên, cậu bị lôi hẳn ra khỏi phòng. ``Tự dưng lại gọi anh làm gì? Anh chỉ là nhân viên lập trình, mấy cái thể thao hay vận động cơ bắp thì anh yếu lắm!''

``Vậy càng tốt!'' nàng Khuyển đẩy cậu ra giữa sân, nhét vào tay anh một chiếc còi nhựa. ``Trọng tài tốt nhất là người không biết gì về môn đó. Mới công bằng!''

``Đúng vậy đó peko!'' nàng Thỏ hưởng ứng, ``Chứ Rushia mà làm trọng tài thì bóng sẽ phát nổ mất.''

``Không phát nổ đâu,'' nàng Pháp sư lẩm bẩm.

Cậu Tổng cũng chỉ biết đường thở dài, tặc lưỡi chấp nhận với đề xuất ``từ trên trời rơi xuống'' này. Trong thâm tâm, cậu thoáng nghĩ: ``\textit{Hy vọng là trận cầu này sẽ không có gì đi quá giới hạn\dots}''

Mt trận đấu bóng chày ba người không mấy hợp lý cũng được sắp đặt. Rushia làm người ném bóng, Pekora làm người đánh, Korone làm người đuổi theo bắt bóng, và Kuon bất đắc dĩ trở thành trọng tài.

Một chiều tháng Mười, không ai biết vì sao họ lại rảnh đến mức làm chuyện này. Có lẽ vì tiết trời quá đẹp, hoặc vì tâm trạng ai đó cần một chút nhẹ nhàng. Cũng có thể chỉ đơn giản là\dots\ hôm ấy, không có lịch ghi hình nào quan trọng.

\ct

``Vị trí nào vào chỗ đó!'' Kuon nói như một người đang bị số phận kéo lê, tay cầm chiếc còi nhựa màu xanh lơ đã mẻ một bên. Cậu đứng bên lề sân, cố giữ khoảng cách an toàn với cả ba người còn lại, không chỉ vì sợ bóng chày đập mặt, mà còn vì một bên là dao găm mang hình hài một nàng chó, một bên là cơn cuồng gào `peko peko', và chính giữa là một tinh linh tử vong trầm cảm. Rõ ràng cậu Tổng không hề tin tưởng rằng đây sẽ là một trận bóng chày bình thường cả.

Ở chính giữa sân, Rushia nắm quả bóng trong tay. Cô hít một hơi thật sâu, mắt nhắm tịt lại như đang tụ linh khí từ cõi u minh. Gió lay nhẹ tà tóc cô, khiến Kuon không nhịn được mà nghĩ: ``\textit{Có lẽ đây là lần đầu trong đời mình thấy ném bóng mà trông giống triệu hồi hồn ma như vậy.}''

``\textbf{Rồi, bắt đầu lượt thứ nhất!}'' cậu hô lên, rồi thổi còi một tiếng.

Rushia lấy đà, rồi ném. Cú ném không mạnh lắm, nhưng đường bóng lại đi lạ kỳ, vòng một đường xoắn nhẹ rồi lướt ngang mặt đất như đang bay lơ lửng.

``Ể\dots\ Ể!?'' cả Pekora và Korone trố mắt trước đường đi không tuân theo logic của quả bóng.

Nàng Thỏ cố vụt mạnh vào quả bóng ngay khi nó tiếp cận, nhưng\dots

*VÚT!*

``Peko!''

\dots cô nàng đập hụt, thậm chí cô còn ngã dập mặt xuống sàn.

``\textbf{Nó không đi thẳng, peko!! Bà ném kiểu gì thế!?}'' cô đứng dậy, phủi bụi trên mặt.

``Làm sao mà tôi biết!'' nàng Pháp sư chớp mắt, tới cô cũng không ngờ cô vừa thực hiện cú ném đó. ``Chắc là tại gió đấy chứ\dots''

Korone từ xa lao tới, lăn một vòng trên đất và tóm được bóng, dù nó chỉ vừa chạm đất nhẹ nhàng. ``\textbf{Gâu gâu! Đuổi bóng là việc của chị mừ\~{}}'' cô hét lớn, bụi đất bay tứ tung.

Cậu Tổng thổi còi một lần nữa, giơ tay lên thông báo: ``Lượt đầu tiên, không tính! Bóng chưa qua strike zone!''

``Strike zone là thế quái nào vậy?'' cả ba cô gái đồng thanh.

``Ờ thì\dots\ chỗ mà người đánh có thể đánh được. Nói tóm lại là đánh trượt!'' Kuon gãi đầu gãi tai, không hề biết gì về khái niệm bóng chày.

``Thì cứ bảo trượt đi, peko!'' nàng Thỏ thốt lên. ``Lại cứ sờ-trai-dôn làm gì\dots''

Lượt đánh thứ hai.

Lần này, Rushia ném nhanh hơn. Bóng bay vèo tới qua một đường parabol hoàn hảo.

``Chuẩn bị tư thế nào\dots'' nàng Thỏ ngước về quả bóng, căn chỉnh chiếc chày.

Nhưng cô chưa kịp vào tư thế đã lùi lại hai bước như thể bị ma ám. ``\textbf{Peko!!}''

Quả bóng nằm gọn trong tay của Korone. ``Sao thế, Peko-chan?''

``N-Nó nhìn em!!'' Pekora chỉ vào vệt đen dưới chân mình. ``Em thề là cái vệt đen đó nhìn em, peko!!''

``Cái gì nhìn em cơ?'' Kuon nhướng mày.

``Cái bóng! Của quả bóng đó, peko!!''

``Do mặt trời đó, đồ ngốc!'' Korone đáp, đang nhai cát như đang tập luyện cho vai chó nhẫn giả trong phim.

``Quả đó cũng tính là trượt nha,'' cậu tuyên bố, lần này dứt khoát hơn.

Lượt đánh thứ ba, cũng là lượt cuối cùng.

Nàng Pháp sư ném cú quyết định. Cô lấy đà chậm rãi, như đang diễn cảnh trong phim thể thao học đường. Khi tay vừa buông bóng, nhạc nền tưởng tượng dường như vang lên, thứ âm thanh bi tráng mà chỉ mình cô nghe thấy.

Bóng bay thẳng, mạnh và mượt.

Pekora hét lên một tiếng xung kích ``\textbf{Peko!!!}'' rồi vung chày thật mạnh.

*CHÁT!*

Cú chạm vang dội đến độ trọng tài đánh rơi cả còi trên tay xuống.

Quả bóng bay lên, xoáy nhẹ trong không trung. Cảnh tượng diễn ra chậm lại trong tâm trí mọi người, ít nhất là theo lời kể của Rushia sau này.

``\textbf{Bóng bay, peko!!}'' Pekora hô.

``\textbf{Bóng bay kìa, gâu gâu!!}'' Korone lao theo, hai chân chạy thoăn thoắt, miệng sủa nhẹ.

Kuon quay lưng lại, giơ tay ôm lấy đầu: ``Đừng đập vào mặt anh là được\dots''

Sau một khoảng khắc kéo dài như thể cả thế kỷ nén lại trong ba giây, Korone bật người lên, với tay.

*PẶP!*

Nàng Khuyển bắt được.

``Yaaay!! Bắt được rồi, bắt được rồi!!'' cô nhảy tưng tưng như chó săn được thưởng bánh quy.

Nàng Pháp sư mỉm cười. Nàng Thỏ đứng sững, gậy vẫn còn giơ trên không. ``Peko\dots\ Em đánh được rồi\dots\''

``Công nhận,'' Kuon vừa nhặt lại chiếc còi, vừa nói với sự vui lây, nhưng rồi giọng cậu bỗng nghiêm túc hẳn, ``cơ mà nếu đây là trận bóng thật, thì em bị loại rồi đó.''

Pekora ngoảnh đầu về phía cậu, trố mắt ngạc nhiên: ``Ể? Nhưng mà em đập được mà?''

``Đúng, nhưng người ta bắt bóng rồi,'' cậu cười chát. ``Ba quả trượt, coi như em bị loại rồi đấy.''

Nàng Thỏ im lặng, như không tin vào những gì cô nghe. ``C-Cái gì chứ, peko\dots?''

``Tiếc quá, Pekora-chan,'' Rushia cười trừ. ``Bà đánh mạnh hơn một chút là được rồi\dots''

Nàng Thỏ cúi đầu, gào lên: ``\textbf{Tại sao chứ, peko!!!???}'' Cô nghiến răng, ngẩng đầu lên nhìn họ đang bụm miệng cười. Hai má cô phồng lên vì tức tối, ném chiếc gậy đánh bóng xuống sàn: ``Khỉ thật chứ, peko! Hai lượt đầu thì hụt, lượt sau thì bị đối phương bắt được!''

Rushia thì vừa lo vừa buồn cười: ``Peko-chan đánh trúng đấy chứ\dots''

Korone vẫn còn đang đeo găng tay, híp mắt: ``Hehe\~{} nhưng mà bị bắt vẫn là bị bắt đó nha\~{}''

``Lần tới em sẽ đánh cho bay luôn khỏi cái văn phòng này cho xem!'' nàng Thỏ xanh tuyên bố, lườm cả hai người còn lại.

Cậu Tổng (kiêm trọng tài) ghi chú vào sổ: ``Lượt thứ ba, Peko-chan đánh trúng. Koro-chan bắt được. Như vậy Ru-chạn vẫn chưa được tính điểm.''

Cậu khẽ liếc đồng hồ, rồi quay sang ba cô gái: ``Được rồi, tới lượt Ru-chan đánh. Ai làm pitcher (pít-chơ, người ném bóng) đây?''
Korone vẫy tay: ``Để chị ném cho!''

\ct

Rushia đứng vào vị trí, tay cầm gậy lóng ngóng như chưa từng chơi bao giờ. Korone thì ở phía đối diện, miệng vẫn cười toe toét, tay vung bóng. Pekora thì lóc cóc đứng sau nàng Khuyển, mặt vẫn chưa hết hậm hực.

``Á! Trượt rồi!'' nàng Pháp sư thốt lên, quay vòng vòng tại chỗ.

Korone ngửa cổ cười: ``Lại lần nữa nha\~{}''

Lần thứ hai, Rushia dốc toàn lực, gồng mình và\dots

*CHÁT!*

Tiếng gậy chạm bóng vang lên. Quả bóng bay lên nhưng xoáy ngang một cách khó hiểu, đập vào tường rồi rơi ngay trước chân Kuon.

Cậu Tổng cúi nhặt lấy, thổi còi thông báo: ``An toàn! Ru-chan được điểm!''

Nàng Pháp sư thở dài, nhưng miệng nở một nụ cười tươi: ``Công nhận chơi bóng chày vui thật đó-nanodesu!''

``Tôi thì chẳng vui chút nào đâu, peko,'' đôi tai thỏ của Pekora cụp xuống, càng tức tối hơn khi cô bạn cùng Thế hệ ghi được điểm.

\ct

Tới lượt Korone.

Nàng Thỏ hăng hái đeo găng, chỉnh tư thế, đứng vào vị trí bắt bóng như một tay chơi chuyên nghiệp. Nàng Pháp sư lùi lại quan sát, còn cậu trọng tài bước ra một bên, chắp tay sau lưng như thể đang xem một màn kịch ngắn.

``Được rồi đó, Pekora-chan\~{} Cẩn thận nha\~{}'' Korone nâng gậy lên, mắt tia tới quả bóng vẫn đang ở trên tay.

*VÚT!*

Pekora ném quả bóng với một lực vừa phải. Rushia chỉ vừa kịp nói ``Đánh nhẹ thôi nha!'' thì\dots

\textbf{*BỐP!!*}

Một cú đánh trời giáng, ngay từ lần đánh đầu tiên. Quả bóng bay khỏi khung hình như một viên đạn.

*XOẢNG!*

Cú đánh bóng đầy uy lực của nàng Khuyển phá tung luôn một ô cửa sổ của văn phòng, tiếp tục hướng ra khỏi khu văn phòng.

``\textbf{Peko!? Mạnh thế!?}'' nàng Thỏ hét lên, lập tức đuổi theo.

``\textbf{Chờ tôi với!}'' Rushia gọi theo, cũng chạy sau Pekora. Korone thấy vậy cũng chạy theo, vừa cười vừa hú lên: ``Khoan khoan, mình đâu có cố ý đánh mạnh thế!''

Cả ba cô gái lao ra khỏi khu quay phim như một cơn lốc ba màu: xanh dương, xanh lục và nâu.
Kuon nhìn theo, đứng im một lát. Anh tháo chiếc mũ trọng tài ra, khẽ thở dài: ``\dots Biết ngay là kiểu gì sẽ bất thường mà.''

Anh đóng sổ lại, quay gót bước chậm về phía máy quay như đang trò chuyện với khán giả (và chính người viết câu chuyện này): ``Cảnh kết thúc. Bản dựng có thể xài được.''

%==========================================TRUYỆN PHỤ=========================================
\trva

Tập truyện ``Cuộc tập dượt bóng chày'' thực chất được lấy cảm hứng từ các khung hình xuất hiện trong tập phim ``Phi vụ 60 giây'' (chính là {\flaregothic ÉPISODE 127}, tập trước đó.)

Theo các kịch bản mà chính Kanata viết cũng như theo suy đoán của người biên dịch, đây là tập phim có thời lượng đặc biệt ngắn, chỉ kéo dài khoảng 60 giây, được biên dịch lại thành kịch bản bởi Lamy, Kanata và Watame trong một buổi tối\dots\ hỗn loạn, với bối cảnh là chính văn phòng \hll. Truyện sau đó được viết lại vào quyển số Mười - hồi Một bởi người biên dịch và sự giúp sức của ChatGPT. Trong đó, người biên dịch sử dụng ba bức hình của tập 127 để phác họa:

\begin{figure}[H]
    \centering
    \begin{subfigure}[t]{0.45\textwidth}
        \centering
        \includegraphics[width=8cm]{../../../../Image/Book?/c-5a.png}
    \end{subfigure}
    \quad
    \begin{subfigure}[t]{0.45\textwidth}
        \centering
        \includegraphics[width=8cm]{../../../../Image/Book?/c-5b.png}
    \end{subfigure}
    \quad
    \begin{subfigure}[t]{0.45\textwidth}
        \centering
        \includegraphics[width=8cm]{../../../../Image/Book?/c-5c.png}
    \end{subfigure}
\end{figure}

Dù ban đầu đây chỉ là một thử nghiệm nhanh, nhưng câu chuyện ngắn xoay quanh Rushia, Korone và Pekora cùng trò bóng chày bất ngờ lại trở thành một trong những bản dựng có triển vọng nhất, và cũng là nguồn cảm hứng cho bản chuyển thể văn học này.

\end{document}